\subsection{The homomorphism F and the shift V}

Let A be a ring. In the rest of this paragraph, we denote respectively by + and × the addition and multiplication laws in W(A). We shall also write 0 for 0 \(_{A}\) and 1 for \(1_{A}\) . We define two maps \(^{1}\) \(F_{A}\) and \(V_{A}\) (also denoted simply by F and V) from W(A) into itself by the formulas

\begin{enumerate}
\def\labelenumi{(\arabic{enumi})}
\setcounter{enumi}{22}
\tightlist
\item
\end{enumerate}

\[
\mathrm{F} _ {\mathbf {A}} (\boldsymbol {a}) = \big (\mathrm{F} _ {n} (a _ {0},..., a _ {n + 1}) \big) _ {n \in \mathbb {N}},\tag{24}
\]

\[
\mathrm{V} _ {\mathrm{A}} (\boldsymbol {a}) = (0, a _ {0}, a _ {1}, \dots)
\]

\((\text{for } \boldsymbol{a} = (a_n)_{n \in \mathbb{N}} \text{ in } \mathrm{W(A)})\) . The map \(V_A\) is called the shift. The formula

\[
\Phi_ {n} \left(\mathrm{F} _ {0} (\boldsymbol {a}), \dots , \mathrm{F} _ {n} (\boldsymbol {a})\right) = \Phi_ {n + 1} \left(a _ {0}, \dots , a _ {n + 1}\right) \quad (n \in \mathbf {N})\tag{25}
\]

follows immediately from (15). It can also be written in the form

\[
\Phi_ {\mathrm{A}} \circ \mathrm{F} _ {\mathrm{A}} = f _ {\mathrm{A}} \circ \Phi_ {\mathrm{A}}.\tag{26}
\]

The formula

\[
\Phi_ {\mathrm{A}} \circ \mathrm{V} _ {\mathrm{A}} = v _ {\mathrm{A}} \circ \Phi_ {\mathrm{A}}\tag{27}
\]

follows from relation (3).

Let \(\rho : \mathbf{B} \to \mathbf{A}\) be a ring homomorphism. The relations

\begin{enumerate}
\def\labelenumi{(\arabic{enumi})}
\setcounter{enumi}{27}
\tightlist
\item
\end{enumerate}

\[
\mathbf {W} (\rho) \circ \mathrm{F} _ {\mathbf {B}} = \mathrm{F} _ {\mathbf {A}} \circ \mathbf {W} (\rho)\tag{29}
\]

\[
\mathrm{W} (\rho) \circ \mathrm{V} _ {\mathrm{B}} = \mathrm{V} _ {\mathrm{A}} \circ \mathrm{W} (\rho)
\]

follow immediately from the definitions.

PROPOSITION 3. — Let A be a ring.

\begin{enumerate}
\def\labelenumi{\alph{enumi})}
\item
  The map \(F_{A}\) is an endomorphism of the ring W(A).
\item
  The map \(V_{A}\) is an endomorphism of the additive group underlying the ring \(\mathbf{W}(\mathbf{A})\) .
\item
  For every a in W(A), one has \(F_{A}(V_{A}(a)) = p.a\) (sum in W(A) of p terms equal to a).
\item
  For all \(a\) and \(b\) in W(A), one has
\end{enumerate}

\begin{enumerate}
\def\labelenumi{(\arabic{enumi})}
\setcounter{enumi}{29}
\tightlist
\item
\end{enumerate}

\[
\mathrm{V} _ {\mathrm{A}} (a \times \mathrm{F} _ {\mathrm{A}} (b)) = \mathrm{V} _ {\mathrm{A}} (a) \times b\tag{31}
\]

\[
\mathrm{V} _ {\mathrm{A}} (a) \times \mathrm{V} _ {\mathrm{A}} (b) = p. \mathrm{V} _ {\mathrm{A}} (a \times b)
\]

(sum in W(A) of p terms equal to \(V_A(a \times b)\)).

\begin{enumerate}
\def\labelenumi{\alph{enumi})}
\setcounter{enumi}{4}
\tightlist
\item
  Set \(\mu = V_{A}(1) = (0, 1, 0, \ldots)\) . For every \(b\) in W(A), one has
\end{enumerate}

\[
\mathrm{V} _ {\mathrm{A}} \left(\mathrm{F} _ {\mathrm{A}} (\boldsymbol {b})\right) = \mu \times \boldsymbol {b}.\tag{32}
\]

\begin{enumerate}
\def\labelenumi{\alph{enumi})}
\setcounter{enumi}{5}
\tightlist
\item
  For every element \(\mathbf{a}\) of W(A), denote by \(\mathbf{a}^{*p}\) the product in W(A) of p elements equal to a. Then one has
\end{enumerate}

\[
\mathrm{F} _ {\mathbf {A}} (\boldsymbol {a}) \equiv \boldsymbol {a} ^ {* p} \bmod p. \mathrm{W} (\mathbf {A}) \quad (\text { ideal of } \mathrm{W} (\mathbf {A}) \text { generated by } p. \mathbf {1}).\tag{33}
\]

Let \(\rho: \mathbf{B} \to \mathbf{A}\) be a ring homomorphism satisfying the conditions of lemma 3 of no. 4. Then \(\mathbf{W}(\rho): \mathbf{W}(\mathbf{B}) \to \mathbf{W}(\mathbf{A})\) is a surjective ring homomorphism, and \(\Phi_{\mathbf{B}}: \mathbf{W}(\mathbf{B}) \to \mathbf{B}^{\mathbf{N}}\) is an injective ring homomorphism. Moreover, \(f_{\mathbf{B}}: \mathbf{B}^{\mathbf{N}} \to \mathbf{B}^{\mathbf{N}}\) is a ring homomorphism. By formulas (26) and (28), one has

\[
\Phi_ {\mathrm{B}} \circ \mathrm{F} _ {\mathrm{B}} = f _ {\mathrm{B}} \circ \Phi_ {\mathrm{B}}, \quad \mathrm{W} (\rho) \circ \mathrm{F} _ {\mathrm{B}} = \mathrm{F} _ {\mathrm{A}} \circ \mathrm{W} (\rho),
\]

whence assertion a) follows immediately. Assertion b) follows analogously from formulas (27) and (29) and from the fact that \(v_{B}\) is an endomorphism of the additive group underlying \(B^{N}\) .

Let \(\pmb{a}\) be an element of W(A), and choose an element \(\pmb{x}\) of W(B) whose image under W(ρ) is \(\pmb{a}\) . Put \(\xi = \Phi_{\mathrm{B}}(\pmb{x})\) . It follows immediately from the definitions of \(f_{\mathrm{B}}\) and \(v_{\mathrm{B}}\) that one has \(f_{\mathrm{B}}(v_{\mathrm{B}}(\xi)) = p.\xi\) (sum in B \(^{\mathbf{N}}\) of \(p\) terms equal to \(\xi\) ). By formulas (26) and (27) (where A is replaced by B), the elements F \(_{\mathrm{B}}\) (V \(_{\mathrm{B}}\) ( \(\pmb{x}\) )) and \(p.\pmb{x}\) of W(B) therefore have the same image \(p.\xi\) under the injective map \(\Phi_{\mathrm{B}}\) , and thus are equal. The formula F \(_{\mathrm{A}}\) (V \(_{\mathrm{A}}\) ( \(\pmb{a}\) )) = \(p.\pmb{a}\) then follows from relations (28) and (29). This proves \(c\) .

Reasoning analogously, one reduces the proof of formula (30) to that of the relation

\[
v _ {\mathrm{B}} (\xi f _ {\mathrm{B}} (\eta)) = v _ {\mathrm{B}} (\xi) \eta
\]

for \(\xi\) , \(\eta\) in B \(^{N}\) . Now this follows from the equalities

\[
\begin{array}{l} \xi f _ {\mathrm{B}} (\eta) = (\xi_ {0} \eta_ {1}, \xi_ {1} \eta_ {2}, \dots) \\ v _ {\mathrm{B}} (\xi) \eta = (0, p \xi_ {0} \eta_ {1}, p \xi_ {1} \eta_ {2}, \dots). \end{array}
\]

Taking (b) and (c) into account, formula (31) follows from formula (30), where \(\boldsymbol{b}\) is replaced by \(\mathrm{V}_{\mathbf{A}}(\boldsymbol{b})\) . Formula (32) is the particular case \(\boldsymbol{a} = \boldsymbol{1}\) of formula (30).

Similarly, one reduces the proof of formula (33) to that of the relation

\[
f _ {\mathrm{B}} (\xi) \equiv \xi^ {p} \bmod p. \Phi_ {\mathrm{B}} (\mathrm{B} ^ {\mathbf {N}}),
\]

where \(\xi^{p}\) denotes the product in \(\mathbf{B}^{\mathbf{N}}\) of \(p\) elements equal to \(\xi\) . By prop. 2, c) of no. 2, this is equivalent to the fact that for every \(n \geqslant 0\) , one has

\[
\sigma \left(\xi_ {n + 1} - \xi_ {n} ^ {p}\right) \equiv \xi_ {n + 2} - \xi_ {n + 1} ^ {p} \bmod p ^ {n + 2} \mathbf {B}.
\]

Now, for every \(n \geqslant 0\) , one has, by loc. cit.,

\[
\sigma (\xi_ {n}) \equiv \xi_ {n + 1} \bmod p ^ {n + 1} \mathbf {B}
\]

since \(\xi = \Phi_{\mathbf{B}}(\boldsymbol{x})\) ; from this one deduces, by means of lemma 1 of no. 1,

\[
\sigma (\xi_ {n}) ^ {p} \equiv \xi_ {n + 1} ^ {p} \bmod p ^ {n + 2} \mathbf {B}.
\]

This proves the desired relation.

Remark. --- For the definition of maps analogous to the maps F and V, in the case of the universal Witt ring, see exercises 36, 37 and 38, p. 52 and following.

